\section{FOL Proofs}

Proofs for FOL will be governed by the same rule system as TFL, with the
addition of some new rules for dealing with singular terms, predicates and
quantifiers. Of course, these rules will take the form of our familiar
introduction and elimination rules, plus some derived rules that let us avoid
repeating frequently needed subproofs.

%% quine: only help yourself to as much structure as you need (p322)

\subsection{Basic Natural Deduction Rules for FOL}

%% TODO: check versus quantifier rules summarized (p315)
%        summary of quantifier rules (p354)

Given the definitions for substitution instances discussed last week, defining
elimination rules for $\eforall x$ and introduction rules for $\eexists x$ is
particularly straight-forward. In our rules, we let \exi be a metavariable for
FOL variables, \etau be a metavariable for FOL terms and \edelta be a
metavariable for FOL dummy names.

With \emph{universal elimination}, we can replace a universally quantified
variable with any term.

\begin{fitchrule}[Universal Elimination ($\eforall$E)]
  \label{def:universal-elimination-rule}
  \leavevmode
  \vspace{-0.9\baselineskip}
  \begin{fitchproof}
    \have[n] {1} {$\eforall\exi\ephi(\exi)$}
    \have[ ] {2} {$\ephi(\etau)$} \Ae{1}
  \end{fitchproof}
\end{fitchrule}

With \emph{existential introduction}, we can replace any one or more instances
of a given name with an existentially-quantified variable.

\begin{fitchrule}[Existential Introduction ($\eexists$I)]
  \label{def:existential-introduction-rule}
  \leavevmode
  \vspace{-0.9\baselineskip}
  \begin{fitchproof}
    \have[n] {1} {$\ephi(\edelta)$}    
    \have[ ] {2} {$\eexists\exi\ephi(\exi)$} \Ee{1}
  \end{fitchproof}
\end{fitchrule}

The two complementary operations, (\emph{existential elimination} and
\emph{universal introduction}), are more constrained in their definitions.

For \emph{universal introduction} ($\eforall$I), we should only allow the rule
to apply when the name in question is \emph{arbitrary}. For the purposes of
proofs, this amounts to not being found in premises or unresolved subproofs
since their use there makes them non-arbitrary.

\begin{fitchrule}[Universal Introduction ($\eforall$I)]
  \label{def:universal-introduction-rule}
  \leavevmode
  \vspace{-0.9\baselineskip}
  \begin{fitchproof}
    \have[n] {1} {$\ephi(\edelta)$}    
    \have[ ] {2} {$\eforall\exi\ephi(\exi)$} \Ee{1}
  \end{fitchproof}
  But only if \edelta does not occur in a premise or an undischarged assumption.
\end{fitchrule}

Finally, with \emph{existential elimination}, if we have \eexists\exi\ephi(\exi)
and can prove \ephi(\edelta)\eif\epsi via subproof, then we can infer \epsi,
provided that it does not contain \edelta.

\begin{fitchrule}[Existential Elimination ($\eexists$E)]
  \label{def:existential-elimination-rule}
  \leavevmode
  \vspace{-0.9\baselineskip}
  \begin{fitchproof}
    \have[n] {1} {$\eexists\exi\ephi(\exi)$}
    \open
    \hypo[i] {2} {$\ephi(\edelta)$} \by{AS}{}
    \have[j] {3} {\epsi}
    \close
    \have[ ] {4} {\epsi} \Ee{1,2-3}
  \end{fitchproof}
  But \edelta must not occur in:
  \begin{enumerate}
  \item any assumption discharged before line $i$
  \item $\eexists\exi\ephi(\exi)$
  \item \epsi
  \end{enumerate}
\end{fitchrule}

\subsection{Subproofs and Dummy Names}

We should pause for a moment to consider the rules for the use of dummy names
(\edelta) in subproofs. Note that both universal introduction and existential
elimination have strict rules about when a dummy name can be used in an
inference step. Let's look at some examples of bad inferences (modified from
\cite{Smi20}).

\begin{fitchproof}
  \hypo      {1} {\eexists xFx} \by{PR}{}
  \open
  \hypo      {2} {Fa}           \by{AS}{}
  \have[**3] {3} {\forall xFx}  \Ai{1,2}
  \close
  \have      {4} {\eforall xFx} \ee{1,2-3}
\end{fitchproof}

We cannot infer that everyone is F just from the fact that someone is. This
proof mistakenly generalizes Fa by misusing the \eforall I rule. The name $a$
is introduced in an assumption that has not yet been discharged (i.e. it is an
assumption of an open subproof.

Here's another one.

\begin{fitchproof}
  \hypo      {1} {\eexists xFx} \by{PR}{}
  \open
  \hypo      {2} {Fa}           \by{AS}{}
  \have      {3} {Fa}           \by{R}{2}
  \close
  \have[**4] {4} {Fa}           \ee{1,2-3}
  \have      {5} {\eforall xFx} \Ai{4}
\end{fitchproof}

This time the problem is that we are trying to draw an instance from a closed
subproof. But the subject of that instance ($a$) is entirely hypothetical,
a representative example introduced by the assumptions of the subproof. The
\eexists E rule does not allow this kind of inference.

Similarly, we cannot reuse representative names. Once we've introduced
$\edelta_i$ via assumption, we need to later use a new name ($\edelta_{i+1}$) if
we want to assume a name for a different existential sentence.

%%% TODO: check line 6 ?
\begin{fitchproof}
  \hypo      {1} {\eexists xFx}           \by{PR}{}
  \hypo      {2} {\eexists xGx}           \by{PR}{}
  \open
  \hypo      {3} {Fa}                     \by{AS}{}
  \open
  \hypo      {4} {Ga}                     \by{AS}{}
  \have      {5} {Fa\eand Ga}             \ai{3,4}
  \have      {6} {}
  \close
  \have[**7] {7} {\eexists x(Fx\eand Gx)} \ee{2,4-6}
  \close
  \have      {8} {\eforall xFx}           \ee{1, 3-7}
\end{fitchproof}

The ``error'' is the reuse of a dummy name in a new subproof (line 4). We cannot
assume that the x who Fs is the same as the x who Gs. But the bad inference step
is actually on line 7. We are not allowed to eliminate dummy names from dummy
sentences when the name is introduced within a closed subproof.

\subsection{Proof Strategies for First-Order Logic}

Similarly, proofs in FOL will proceed with a strategy based on the one we
devised for FOL, with the addition of some new techniques to deal with
quantifiers. Quite tersely, we can describe these new techniques as follows.

\begin{enumerate}
\item Work forward when dealing with the existential quantifier (\eexists)
\item Work backward when dealing with universal quantifiers (\eforall)
\end{enumerate}

Consider this example proving $\eforall x(Px\eif Q)\eproves\eexists x Px\eif Q$.

First we set up the proof with premise and conclusion at their respective ends.

\begin{fitchproof}
  \hypo      {1} {$\eforall x(Px\eif Q)$} \by{PR}{}
  \have[ ]   {2} {$\vdots$}
  \have[n]   {3} {$\eexists x Px\eif Q$}
\end{fitchproof}

Then we can work backward from the conclusion, a conditional.

\begin{fitchproof}
  \hypo      {1} {$\eforall x(Px\eif Q)$} \by{PR}{}
  \open
  \hypo      {2} {$\eexists x Px$}          \by{AS}{}
  \have[ ]   {3} {$\vdots$}
  \have[m]   {4} {$Q$}
  \close
  \have[n]   {5} {$\eexists x Px\eif Q$} \ie{2-4}
\end{fitchproof}

Now our best move is to work forward from the existential in the subproof,
adding a new subproof with an instance as its assumption and $Q$ as its
conclusion.

\begin{fitchproof}
  \hypo      {1} {$\eforall x(Px\eif Q)$} \by{PR}{}
  \open
  \hypo      {2} {$\eexists x Px$}          \by{AS}{}
  \open
  \hypo      {3} {$Pa$}                     \by{AS}{}
  \have[ ]   {4} {$\vdots$}
  \have[l]   {5} {$Q$}
  \close
  \have[m]   {6} {$Q$}                     \ee{2,3-5}
  \close
  \have[n]   {7} {$\eexists x Px\eif Q$} \ie{2-}
\end{fitchproof}

Now we can unpack the universal quantifier inside the subproof, working forward
from the first line. We can use $a$ for its instance, which will be helpful,
since the conclusion follows by modus ponens.

\begin{fitchproof}
  \hypo      {1} {$\eforall x(Px\eif Q)$} \by{PR}{}
  \open
  \hypo      {2} {$\eexists x Px$}          \by{AS}{}
  \open
  \hypo      {3} {$Pa$}                     \by{AS}{}
  \have      {4} {$Pa\eif Q$}               \Ae{1}
  \have      {5} {$Q$}                      \ie{3,4}
  \close
  \have      {6} {$Q$}                     \ee{2,3-5}
  \close
  \have      {7} {$\eexists x Px\eif Q$} \ie{2-6}
\end{fitchproof}

After merging the new proof strategies into the ones for TFL in Obs.
\ref{obs:tfl-proof-strategies}, we get the following complete list of proof
strategies for first-order logic.

\begin{observation}[Proof Strategies for FOL]
  \label{obs:fol-proof-strategies}
  We can tersely describe the basic strategy for FOL proofs as follows.
  \begin{outline}[enumerate]

    \1 Work backwards from the conclusion.
    \2 Conjunction: what components needed to use the \eand I rule as the
    last step? i.e. what are the conjuncts?
    \2 Conditional: what are the assumption and conclusion of the subproof
    we need to take the last step? i.e. what are the antecedant and consequent
    of the conditional?
    \2 Negation: what subproof is needed to show a contradiction?
    \2 Disjunction: which of the disjuncts looks easiest to prove?
    \2 Universal quantification: how can we prove an arbitrary instance?
    
    \1 Work forwards from the current step.
    \2 Conjunction: unpack the conjuncts
    \2 Disjunction: what sentence do you need to prove from both disjuncts?
    \2 Conditional: how can you prove the antecedant?
    \2 Negation: how can you prove a contradiction?
    \2 Existential quantification: how can we work from an assumed instance to
    the conclusion?
    
    \1 Leverage \ebot
    \2 Work toward a contradiction to conclude anything, especially in a
    subproof
    \2 Indirect proof: build a contradictory subproof that assumes the negation

  \end{outline}
\end{observation}

\subsection{Some Examples of FOL Proofs}

Let's looks at a few more examples of proofs in first-order logic.

First, we'll prove $\eexists xFx,\eforall x(Fx\eif Gx)\eproves\eexists xGx$.

\begin{fitchproof}
  \hypo      {1} {$\eexists xFx$}          \by{PR}{}
  \hypo      {2} {$\eforall x(Fx\eif Gx)$} \by{PR}{}
  \have      {3} {$\vdots$}
  \have[n]   {4} {$\eexists xGx$}
\end{fitchproof}

Working forward from the first line, we start a subproof assuming $Fa$. And
working backward from the last line, we end that subproof with $\eexists xGx$.

\begin{fitchproof}
  \hypo      {1} {$\eexists xFx$}          \by{PR}{}
  \hypo      {2} {$\eforall x(Fx\eif Gx)$} \by{PR}{}
  \open
  \hypo      {3} {$Fa$}                    \by{AS}{}
  \have      {4} {$\vdots$}
  \have[m]   {5} {$\eexists xGx$}
  \close
  \have[n]   {6} {$\eexists xGx$}          \ee{1,3-5}
\end{fitchproof}

Now we can unpack the universal with an instance using $a$, get $Ga$ by modus
ponens. The existential generalization $\eexists xGx$ follows immediately. And
we are done.

\begin{fitchproof}
  \hypo      {1} {$\eexists xFx$}          \by{PR}{}
  \hypo      {2} {$\eforall x(Fx\eif Gx)$} \by{PR}{}
  \open
  \hypo      {3} {$Fa$}                    \by{AS}{}
  \have      {4} {$Fa\eif Ga$}             \Ae{2}
  \have      {5} {$Ga$}                    \ie{3,4}
  \have      {6} {$\eexists xGx$}          \ei{5}
  \close
  \have      {7} {$\eexists xGx$}          \ee{1,3-6}
\end{fitchproof}

We can show the validity of $\eforall x(Hx\eif\enot Gx),\forall x(Fx\eif Gx)
\eproves \eforall x(Hx\eif\enot Fx )$ with a more involved proof.

\begin{fitchproof}
  \hypo      {1} {$\eforall x(Hx\eif\enot Gx)$}   \by{PR}{}
  \hypo      {2} {$\eforall x(Fx\eif Gx)$}         \by{PR}{}
  \have      {3} {$\vdots$}
  \have[n]   {4} {$\eforall x(Hx\eif\enot Fx )$}
\end{fitchproof}

Working backward from the conclusion suggests that we are trying to prove an
arbitrary instance $(Ha\eif\enot Fa)$.

\begin{fitchproof}
  \hypo      {1} {$\eforall x(Hx\eif\enot Gx)$}   \by{PR}{}
  \hypo      {2} {$\eforall x(Fx\eif Gx)$}         \by{PR}{}
  \have      {3} {$\vdots$}
  \have[m]   {4} {$(Ha\eif\enot Fa )$}
  \have[n]   {5} {$\eforall x(Hx\eif\enot Fx )$}  \Ai{4}
\end{fitchproof}

Let's unpack the premises with that instance in mind.

\begin{fitchproof}
  \hypo      {1} {$\eforall x(Hx\eif\enot Gx)$}   \by{PR}{}
  \hypo      {2} {$\eforall x(Fx\eif Gx)$}         \by{PR}{}
  \have      {3} {$Ha\eif\enot Gx$}               \Ae{1}
  \have      {4} {$Fa\eif Ga$}                    \Ae{2}
  \have      {5} {$\vdots$}
  \have[m]   {6} {$(Ha\eif\enot Fa )$}
  \have[n]   {7} {$\eforall x(Hx\eif\enot Fx )$}  \Ai{6}
\end{fitchproof}

Now we can build a subproof for $Ha\eif\enot Fa$.

\begin{fitchproof}
  \hypo      {1} {$\eforall x(Hx\eif\enot Gx)$}   \by{PR}{}
  \hypo      {2} {$\eforall x(Fx\eif Gx)$}         \by{PR}{}
  \have      {3} {$Ha\eif\enot Ga$}               \Ae{1}
  \have      {4} {$Fa\eif Ga$}                    \Ae{2}
  \open
  \hypo      {5} {$Ha$}                           \by{AS}{}
  \have      {6} {$\vdots$}
  \have[l]   {7} {$\enot Fa$}
  \close
  \have[m]   {8} {$(Ha\eif\enot Fa )$}            \ie{5-7}
  \have[n]   {9} {$\eforall x(Hx\eif\enot Fx )$}  \Ai{8}
\end{fitchproof}

Now we will use modus ponens and modus tollens with those conditionals we
unpacked before. And we are done.

\begin{fitchproof}
  \hypo      {1} {$\eforall x(Hx\eif\enot Gx)$}   \by{PR}{}
  \hypo      {2} {$\eforall x(Fx\eif Gx)$}         \by{PR}{}
  \have      {3} {$Ha\eif\enot Ga$}               \Ae{1}
  \have      {4} {$Fa\eif Ga$}                    \Ae{2}
  \open
  \hypo      {5} {$Ha$}                           \by{AS}{}
  \have      {6} {$\enot Ga$}                     \ie{3,5}
  \have      {7} {$\enot Fa$}                     \by{MT}{4,6}
  \close
  \have      {8} {$(Ha\eif\enot Fa )$}            \ii{5-7}
  \have      {9} {$\eforall x(Hx\eif\enot Fx )$}  \Ai{8}
\end{fitchproof}

Things get more difficult when we start to apply more subproof-oriented
strategies. Consider this example from \cite[][p. 308]{Smi20}.

\begin{fitchproof}
  \hypo {1} {\eexists xFx\eor\exists xGx} \by{PR}{}
  \open
  \hypo {2} {\eexists xFx}                \by{AS}{}
  \open
  \hypo {3} {Fa}                          \by{AS}{}
  \have {4} {Fa\eor Ga}                   \oi{3}
  \have {5} {\eexists x(Fx\eor Gx)}       \Ei{4}
  \close
  \have {6} {\eexists x(Fx\eor Gx)}       \Ee{2,3-5}
  \close
  \open
  \hypo {7} {\eexists xgx}                \by{AS}{}
  \open
  \hypo {8} {Ga}                          \by{AS}{}
  \have {9} {Fa\eor Ga}                   \oi{6}
  \have {10}{\eexists x(Fx\eor Gx)}       \Ei{7}
  \close
  \have {11} {\eexists x(Fx\eor Gx)}      \Ee{7,8-10}
  \close
  \have {12} {\eexists x(Fx\eor Gx)}      \oe{1,2-6,7-11}
\end{fitchproof}

\subsection{Generalizing from Proofs}

As with TFL, we can prove the validity of entire classes of arguments in FOL
by using \emph{substitution instances}. Indeed, by using metalogical notation in
place of FOL formulas, we can generalize our arguments into \emph{schemas}.

\begin{observation}[FOL wffs can be generalized into metalogical schemas]
  Since the validity of a formal argument is dependent on its shape alone (i.e.
  its form), any valid sentence (or argument) can be reformulated as a
  \emph{schema} by consistently replacing expressions in the sentence with
  metalogical variables of the appropriate type.
\end{observation}

One corollary of this observation is that we can use the same practice with
proofs.

\begin{observation}[Substitution Instances of Proofs are Proofs]
  For any proof \eGamma in the natural deduction system of first-order logic, a
  substitution instance $\eGamma'$ of \eGamma is also a proof.
\end{observation}

As a result, we can just prove the schema itself. Given all the different ways
we can do a consistent substitution, this means we can often prove arguments for
formulas that are superficially quite distinct. Consider this example from
\cite{Smi20}, where the justifications are omitted for space, but should be
obvious enough.

\threeitems{
    \begin{fitchproof}
      \hypo {1} {\eforall xFx}
      \hypo {2} {\eforall x(Fx\eif Gx)}
      \have {3} {Fa}
      \have {4} {Fa\eif Ga}
      \have {5} {Ga}
      \have {6} {\eforall xGx}
    \end{fitchproof}
}{
\begin{fitchproof}
  \hypo {1} {\eforall\exi\ephi(\exi)}
  \hypo {2} {\eforall\exi(\ephi(\exi)\eif\epsi(\exi))}
  \have {3} {\ephi(\etau)}
  \have {4} {\ephi(\etau)\eif\epsi(\etau)}
  \have {5} {\epsi(\etau)}
  \have {6} {\eforall\exi\epsi(\exi)}
\end{fitchproof}
}{
\begin{fitchproof}
  \hypo {1} {\eforall z(\eexists yRzy\eand Hz)}
  \hypo {2} {\eforall z((\eexists yRzy\eand Hz)\eif\eforall xSxz)}
  \have {3} {\eexists yRcy\eand Hc}
  \have {4} {(\eexists yRcy\eand Hc)\eif\eforall xSxc}
  \have {5} {\eforall xSxc}
  \have {6} {\eforall z\eforall xSxz}
\end{fitchproof}

}

We will call such a schematic proof a \emph{proof schema}. Just like proofs
in the object language, a proof using proof schemas shows the validity of the
arguments that they prove. As in TFL, we can convert back and forth between
schematic theorems and theorems in the object language. The advantage of the
schematic proofs is that they distill the problem as much as possible to a mere
matter of form. The relationships that matter are revealed as directly as
possible.

\subsection{Equivalences between Quantified Sentences}

Like conjunction and disjunction, existential and universal quantifiers are
mutually-definable with the help of the negation operator. We saw the conversion
rules for quantifers in Observation \ref{obs:quantifer-duals}, but we did not
prove them. But we can do so now with natrual deduction. As a reminder:

\begin{observation}[Duals of Negative Quantified Sentences]
  \label{obs:quantifer-duals-formal}
  The quantified sentences of FOL can be related via negation with the following
  equivalences:
  \begin{enumerate}
  \item $\enot\eforall\exi\ephi\eiff\eexists\exi\enot\ephi$
  \item $\enot\eexists\exi\ephi\eiff\eforall\exi\enot\ephi$
  \end{enumerate}
\end{observation}

As usual, we will prove these results one conditional at a time. We will
express the sentences schematically, but we could easily do the same proof
with $Fx$ and $Fa$ in place of \ephi(\exi) and \ephi(\edelta), respectively.
\footnote{Similalrly for \ephi(\exi) and \ephi(\etau).}

\subsubsection{$\enot\eforall\exi\ephi\eiff\eexists\exi\enot\ephi$}

\twoitems{
  \begin{fitchproof}
    \hypo   {1} {\enot\eforall\exi\ephi(\exi)}      \by{PR}{}
    \open
    \hypo   {2} {\enot\eexists\exi\enot\ephi(\exi)} \by{AS}{}
    \open
    \hypo   {3} {\enot\ephi(\etau)}                 \by{AS}{}
    \have   {4} {\eexists\exi\enot\ephi(\exi)}      \ei{3}
    \have   {5} {\ebot}                             \ne{2,4}
    \close
    \have   {6} {\enot\enot\ephi(\etau)}            \by{IP}{3-5}
    \have   {7} {\eforall\exi\ephi(\exi)}            \Ai{6}
    \have   {8} {\ebot}                             \ne{1,7}
    \close
    \have   {9} {\eexists\exi\enot\ephi(\exi)}      \by{IP}{2-8}
  \end{fitchproof}
}{
  \begin{fitchproof}
    \hypo   {1}  {\eexists\exi\enot\ephi(\exi)}  \by{PR}{}
    \open
    \hypo   {2}  {\eforall\exi\ephi(\exi)}        \by{AS}{}
    \open
    \hypo   {3}  {\enot\ephi(\edelta)}             \by{AS}{}
    \have   {4}  {\ephi(\edelta)}                  \Ae{2}
    \have   {5}  {\ebot}                           \ne{3,4}
    \close
    \have   {6}  {\ebot}                           \ee{1,3-5}
    \close
    \have   {7} {$\enot\eforall\exi\ephi$}         \ni{2-6}
  \end{fitchproof}
}

\subsubsection{$\enot\eexists\exi\ephi\eiff\eforall\exi\enot\ephi$}
  
\twoitems{
  \begin{fitchproof}
    \hypo {1} {\enot\eexists\exi\ephi}        \by{PR}{}
    \open
    \hypo {2} {\ephi(\etau)}                  \by{AS}{}
    \have {3} {\eexists\exi\ephi(\etau)}      \ei{2}
    \have {4} {\ebot}                         \ne{1,3}
    \close
    \have {5} {\not\ephi(\etau)}              \ni{2-4}
    \have {6} {\eforall\exi\enot\ephi(\exi)}  \Ai{5}
  \end{fitchproof}
}{
  \begin{fitchproof}
    \hypo {1} {\eforall\exi\enot\ephi(\exi)}   \by{PR}{}
    \open
    \hypo {2} {\eexists\exi\ephi(\exi)}        \by{AS}{}
    \open
    \hypo {3} {\ephi(\edelta)}                 \by{AS}{}
    \have {4} {\enot\ephi(\edelta)}            \Ae{1}
    \have {5} {\ebot}                          \ne{3,4}
    \close
    \have {6} {\ebot}                          \ee{2,3-5}
    \close
    \have {7} {$\enot\eexists\exi\ephi(\exi)$} \ni{2-6}
  \end{fitchproof}
}

\subsection{Quantifier Conversion Rule}

Now that we've proved the schemas valid, we can introduce natural deduction
rules for the conversion of quantifiers as follows.

\begin{fitchrule}[Conversion of Quantifiers ($CQ$))]
  \label{def:quantifier-conversion-rule}
  We can relate the quantifiers of first-order logic, with the following two
  pairs of complementary rules. One pair for handling a negated existential
  quantifier:\\
  \leavevmode
  %\vspace{-0.9\baselineskip}
  \twoitems{
    \begin{fitchproof}
      \have[n] {1} {$\eforall\exi\enot\ephi$}
      \have[ ] {2} {$\enot\eexists\exi\ephi$} \by{CQ}{1}
    \end{fitchproof}
  }{
    \begin{fitchproof}
      \have[n] {1} {$\enot\eexists\exi\ephi$}
      \have[ ] {2} {$\eforall\exi\enot\ephi$} \by{CQ}{1}
    \end{fitchproof}
  }
  \\
  and one pair for handling a negated universal quantifier: \\
  \twoitems{
    \begin{fitchproof}
      \have[n] {1} {$\eexists\exi\enot\ephi$}
      \have[ ] {2} {$\enot\eforall\exi\ephi$} \by{CQ}{1}
    \end{fitchproof}
  }{
    \begin{fitchproof}
      \have[n] {1} {$\enot\eforall\exi\ephi$}
      \have[ ] {2} {$\eexists\exi\enot\ephi$} \by{CQ}{1}
    \end{fitchproof}
  }
\end{fitchrule}

We will use the rules in a few examples. First, since we use the two kinds of
quantifiers differently in quantified sentences, we'll prove
$\enot\eexists x(Fx\eand Gx)\eproves\eforall x(Fx\eif Gx)$.

\begin{fitchproof}
  \hypo {1} {\enot\eexists x(Fx\eand Gx)} \by{PR}{}
  \have {2} {\eforall x\enot(Fx\eand Gx)} \by{CQ}{1}
  \open
  \hypo {3} {Fa}                          \by{AS}{}
  \open
  \hypo {4} {Ga}                          \by{AS}{}
  \have {5} {Fa\eand Ga}                  \ai{3,4}
  \have {6} {\enot(Fa\eand Ga)}           \Ae{2}
  \have {7} {\ebot}                       \ne{5,6}
  \close
  \have {8} {\enot Ga}                    \ni{4-7}
  \close
  \have {9} {Fa\eif Ga}                   \ii{3-8}
  \have {10}{\eforall x(Fx\eif Gx)}       \Ai{9}
\end{fitchproof}

Similarly, working with the contents of a complex expression that is within the
scope of a quantifier is challenging. Quantifier conversion rules make it easier
to get them to interact in a way that aligns better with the natural deduction
rules. Consider this proof of
$\eforall x(Fx\eor Gx),\enot\eforall xFx\eproves\eexists xGx$ where line 2 is
inert until we apply CQ.

\begin{fitchproof}
  \hypo {1} {\eforall x(Fx\eor Gx)} \by{PR}{}
  \hypo {2} {\enot\eforall xFx}     \by{PR}{}
  \have {3} {\eexists x\enot Fx}    \by{CQ}{2}
  \open
  \hypo {4} {\enot Fa}              \by{AS}{}
  \have {5} {Fa\eor Ga}             \Ae{1}
  \have {6} {Ga}                    \by{DS}{}
  \have {7} {\eexists xGx}          \Ei{6}
  \close
  \have {8} {\eexists xGx}          \ee{3,4-7}
\end{fitchproof}

We can even use CQ twice to reach inside universally quantified sentences in a
way reminicent of modus tollens:
$\eforall x(Fx\eif Gx),\enot\eforall xGx\eproves\enot\eforall xFx$

\begin{fitchproof}
  \hypo {1} {\eforall x(Fx\eif Gx)} \by{PR}{}
  \hypo {2} {\enot\eforall xGx}     \by{PR}{}
  \have {3} {\eexists x\enot Gx}    \by{CQ}{}
  \open
  \hypo {4} {\enot Ga}              \by{AS}{}
  \have {5} {Fa\eif Ga}             \Ae{1}
  \have {6} {\enot Fa}              \by{MT}{4,5}
  \have {7} {\eexists x\enot Fx}    \Ei{6}
  \close
  \have {8} {\eexists xFx}          \Ee{3,4-7}
  \have {9} {\enot\eforall xFx}     \by{CQ}{}
\end{fitchproof}

Derived rules can simplify to triviality a proof that would be extremely
convoluted with just the basic rules. Consider
$\enot\eforall x\enot Fx\eproves\eexists xFx$.

\begin{fitchproof}
  \hypo {1} {\enot\eforall x\enot Fx} \by{PR}{}
  \have {2} {\eexists x\enot\enot Fx} \by{CQ}{}
  \open
  \hypo {3} {\enot\enot Fa}           \by{AS}{}
  \have {4} {Fa}                      \by{DNE}{}
  \have {5} {\eexists xFx}            \Ei{4}
  \close
  \have {} {\eexists xFx}             \Ee{2,3-5}
\end{fitchproof}

\subsection{Proofs of Inferences with Quantified Sentences}

The natutal deduction system makes it possible to prove the validity of a
variety of inferences involving quantified sentences. First, we can show that
a universal generalization proves the related existential generalization.
$\eforall xFx\eproves\eexists xFx$, which is the expected result for a non-empty
domain. The proof is a trivial application of the rules.

\begin{fitchproof}
  \hypo {1} {\eforall xFx} \by{PR}{}
  \hypo {2} {Fa}           \Ae{1}
  \have {3} {\eexists xFx} \Ei{2}
\end{fitchproof}

The system starts to get more complex with multiple quantifiers. But the simple
cases are still pretty easy to prove. For example:
$\eexists y\eforall xLxy\eproves\eforall x\eexists yLxy$

\begin{fitchproof}
  \have {1} {\eexists y\eforall xLxy} \by{PR}{}
  \open
  \hypo {2} {\eforall xLxa}           \by{AS}{}
  \have {3} {Lba}                     \Ae{2}
  \have {4} {\eexists yLby}           \Ei{3}
  \have {5} {\eforall x\eexists yLxy} \Ai{4}
  \close
  \hypo {6} {\eforall x\eexists yLxy} \Ee{1,2-5}
\end{fitchproof}

The reverse is invalid, as we proved earlier with a counterexample (see: Fig.
\ref{mod:minimal-symmetric}).

Finally, consider this example which shows how a universal implication can be
factored out of a disjunction, simplifying the formula significantly.
$\eforall x(Fx\eif Gx)\eor\eforall x(Fx\eif Hx)\eproves
\eforall x(Fx\eif(Gx\eor Hx))$

\begin{fitchproof}
  \hypo {1} {\eforall x(Fx\eif Gx)\eor\eforall x(Fx\eif Hx)} \by{PR}{}
  \open
  \hypo {2} {\eforall x(Fx\eif Gx)}                          \by{AS}{}
  \have {3} {Fa \eif Ga}                                     \Ae{2}
  \open
  \hypo {4} {Fa}                                             \by{AS}{}
  \have {5} {Ga}                                             \ie{3,4}
  \have {6} {Ga\eor Ha}                                      \oi{5}
  \close
  \have {7} {Fa \eif (Ga\eor Ha)}                            \ii{4-6}  
  \have {8} {\eforall x(Fx\eif(Gx\eor Hx))}                  \Ai{7}
  \close
  \open
  \hypo {9} {\eforall x(Fx\eif Hx)}                          \by{AS}{}
  \have {10}{Fa \eif Ga}                                     \Ae{2}
  \open
  \hypo {11}{Fa}                                             \by{AS}{}
  \have {12}{Ha}                                             \ie{10,11}
  \have {13}{Ga\eor Ha}                                      \oi{12}
  \close
  \have {14}{Fa \eif (Ga\eor Ha)}                            \ii{11-13}  
  \have {15}{\eforall x(Fx\eif(Gx\eor Hx))}                  \Ai{7}
  \close
  \have {16} {\eforall x(Fx\eif(Gx\eor Hx))}                 \by{DS}{1,2-8,9-15}
\end{fitchproof}

\subsection{Identity}

Finally, we will need some rules for handling the identity predicate ($=$). As
usual, these will take to form of introduction and elimination rules for $=$.

If \etau is a term, we can add a sentence at any point that says $\etau=\etau$.

\begin{fitchrule}[Identity Introduction ($=$I)]
  \label{def:identity-introduction-rule}
  \leavevmode
  \vspace{-0.9\baselineskip}
  \begin{fitchproof}
    \have[ ] {1} {$\etau=\etau$} \by{$=$I}{}
  \end{fitchproof}
\end{fitchrule}

And, as with existential introduction, we can perform \emph{identity
elimination} on any one or more instances of the relevant term. Whenever we
know that $\etau_1=\etau_2$ for terms $\etau_1$ and $\etau_2$, we can replace
all of the occurences of $\etau_1$ in $\ephi(\etau_1)$ with $\etau_2$ and the
resulting $\ephi(\etau_2)$ will have exactly the same truth value. $=$E is also
known as \emph{Leibniz's Law}.

\begin{fitchrule}[Leibniz's Law ($=$E)]
  \label{def:identity-elimination-rule}
  Since identity is a symmetric relation, we define a pair of rules for the
  elimination of non-trivial identity sentences like $\etau=\edelta$. \\
  \leavevmode
  \vspace{-0.9\baselineskip}
  \twoitems{
    \begin{fitchproof}
      \have[m] {1} {$\etau=\edelta$}
      \have[n] {2} {$\ephi(\etau)$}    
      \have[ ] {3} {$\ephi(\edelta)$} \by{$=$E}{1,2}
    \end{fitchproof}
  }{
    \begin{fitchproof}
      \have[m] {1} {$\etau=\edelta$}
      \have[n] {2} {$\ephi(\edelta)$}    
      \have[ ] {3} {$\ephi(\etau)$}   \by{$=$E}{1,2}
    \end{fitchproof}
  }
\end{fitchrule}

As always, the best way to understand the rules is to see them used in proofs.

Here's a trivial one that we could use to prove identity is symmetric:
$a=b\eproves b=a$. Recall that $=$ is just a two-place predicate, so line 2 is
playing \textbf{that} role in the inference to line 3.

\begin{fitchproof}
  \hypo {1} {a=b} \by{PR}
  \have {2} {a=a} \by{$=$I}
  \have {3} {b=a} \by{$=$E}{1,2}
\end{fitchproof}

We can also prove that $=$ is transitive. Here we will prove the generalization,
not just an arbitrary instance.

\begin{fitchproof}
  \open
  \hypo {1} {$a=b \eand b=c$}                                 \by{AS}{}
  \have {2} {$a=b$}                                           \ae{1}
  \have {3} {$b=c$}                                           \ae{1}
  \have {4} {$a=c$}                                           \by{$=$E}{2,3}
  \close
  \have {5} {$(a=b \eand b=c)\eif a=c$}                       \ii{1-4}
  \have {6} {$\eforall z((a=b \eand b=c)\eif a=c)$}           \Ai{5}
  \have {7} {$\eforall y\eforall z((a=b \eand b=c)\eif a=c)$} \Ai{6}
  \have {8} {$\eforall x\eforall y\eforall z((a=b \eand b=c)\eif a=c)$} \Ai{7}
\end{fitchproof}

So we can say that
$\eproves\eforall x\eforall y\eforall z((a=b \eand b=c)\eif a=c)$ is a
\emph{theorem} of first-order logic.

\begin{definition}[Theorem of First-Order Logic]
  A \emph{theorem} in first-order logic is any FOL sentence \ephi which can be
  proven without premises (i.e. $\eproves\ephi$).
\end{definition}

Not all theorems are as straightforward to prove as the symmetry and
transitivity of identity, however. Many proceed inirectly and involve multiple
layers of subproof. For example, consider this proof of
$\eproves\enot\eexists x\eforall y((Lxy\eif\enot Lyy)\eand(\enot Lyy\eif Lxy))$,
which is short but far from simple.

\begin{fitchproof}
  \open
  \hypo {1} {\eexists x\eforall y(
    (Lxy\eif\enot Lyy)\eand(\enot Lyy\eif Lxy))}                   \by{AS}{}
  \open
  \hypo {2} {\eforall y((Lay\eif\enot Lyy)\eand(\enot Lyy\eif Lay))} \by{AS}{}
  \have {3} {(Laa\eif\enot Laa)\eand(\enot Laa\eif Laa)}             \Ae{2}
  \have {4} {Laa\eif\enot Laa}                                       \ae{3}
  \have {5} {\enot Laa\eif Laa}                                      \ae{3}
  \open
  \hypo {6} {Laa}                                                    \by{AS}{}
  \have {7} {\enot Laa}                                              \ie{4,6}
  \have {8} {\ebot}                                                  \ne{6,7}
  \close
  \have {9} {\enot Laa}                                              \ni{6-8}
  \have {10}{Laa}                                                    \ie{5,9}
  \have {11}{\ebot}                                                  \ne{9,10}
  \close
  \have {12}{\ebot}                                                  \Ee{1,2-11}
  \close
  \have {13}{\enot\eexists x\eforall y(
    (Lxy\eif\enot Lyy)\eand(\enot Lyy\eif Lxy))}                   \by{IP}{1-12}

\end{fitchproof}

\subsection{Basic FOL Metalogic}

Without getting too deep into the metalogic, we will make some basic claims
about first-order logic.

First, note that the system never fails to provide a truth value for well-formed
formulas of the language. Whenever we have a sentence \ephi built by the grammar
of first-order logic, we know that $\elib\ephi\erib\ein\mathbf{T}$. Consider
this informal proof adapted from \cite[][pp. 342-343]{Smi20}.

\begin{observation}[Every sentence of FOL has a truth value]
  Let $\mathcal(M)=\langle\boldsymbol{T},\boldsymbol{D},\elib\cdot\erib\rangle$
  be a model for FOL $\mathcal{L}$. We will say that $\ephi\ein\mathcal{L}$ just
  in case \ephi is a well-formed formula of $\mathcal{L}$. We will say that
  $\elib\ephi\erib\ein\mathbf{T}$ just in case \ephi is assigned a truth value
  by $\elib\cdot\erib$.

  Now suppose we have a wff \epsi. We know that we can iteratively break \epsi
  down into a parse tree $\mathcal{T}$ by reversing the rule associated with its
  main operator \etheta. We also know that the leaves of $\mathcal{T}$ will each
  be an atomic sentence $\echi(\edelta_i)$ with dummy names replacing all the
  variables. We further know that $\elib\echi(\edelta_i)\erib^{[n/\delta_i]}$ is
  defined for every \echi that can be built by the grammar of $\mathcal{L}$,
  since it has been updated to account for the dummy names $\edelta_i$.

  Reversing our course in the parse tree, we know by the rules of FOL semantics
  that we can compute the truth of a complex sentence built by FOL connectives
  on the basis of the truth of the sentences that are connected. Since we know
  the truth value at each of the leaves, computing the truth of any node is just
  a matter of working up the part of the tree that it \emph{dominates}. When we
  reach the top of the tree, the truth value that we compute there will be the
  truth value of the entire formula.
\end{observation}

A more complete proof would walk through the rules in detail: how do each of
the rules for a parse tree preserve the FOL syntax? and how do each of the rules
for computing FOL truth preserve the FOL semantics? But the informal sketch
gives us a good idea how that might proceed.

\begin{definition}[Validity and Entailment]
  Let $\mathcal{L}$ be a system of first-order logic. Also let \epsi be a
  sentence of $\mathcal{L}$ and $\eGamma=\{\ephi_1,\dots,\ephi_n\}$ be a set of
  sentences of $\mathcal{L}$. The inference $\eGamma\therefore\epsi$ is
  \emph{valid} iff there is no model
  $\mathcal(M)=\langle\mathbf{T},\mathbf{D},\elib\cdot\erib\rangle$
  where $\elib\ephi_i\erib=T$ for all $\ephi_i\ein\eGamma$ but
  $\elib\epsi\erib=F$.

  When $\eGamma\therefore\epsi$ is valid, we say that \eGamma \emph{entails}
  \epsi ($\eGamma\eentails\epsi$).
\end{definition}

Similarly, we can define a notions like \emph{tautology} and \emph{consistency}
against a backdrop of possible models.

\begin{definition}[Tautology]
  \label{def:fol-logical-truth}
  We will say that an FOL sentence \ephi is a \emph{tautology} iff
  $\elib\ephi\erib=T$ on \textbf{every} possible model
  $\mathcal(M)=\langle\mathbf{T},\mathbf{D},\elib\cdot\erib\rangle$.
\end{definition}

\begin{definition}[Consistency]
  We will say that a set of FOL sentences $\eGamma=\{\ephi_1,\dots,\ephi_n\}$ is
  \emph{consistent} iff
  $\elib\ephi_i\erib=T$ for all $\ephi_i\ein\eGamma$ on \textbf{some} possible
  model $\mathcal(M)=\langle\mathbf{T},\mathbf{D},\elib\cdot\erib\rangle$.
\end{definition}

Finally note that FOL is both \emph{sound} and \emph{complete}.

\begin{observation}[Soundness and Completeness of First-Order Logic]
  The system comprising the syntax and semantics of first-order logic is both
  \emph{sound} and \emph{complete}
  \begin{argument}
  \item[ ] $ephi_1,\dots,\ephi_n\eproves\epsi$ iff
    $ephi_1,\dots,\ephi_n\eentails\epsi$
  \end{argument}
  The observation consists of two parts. First, our system is \emph{sound} in
  that every provable argument is also a valid argument.
  \begin{argument}
  \item[ ] if $ephi_1,\dots,\ephi_n\eproves\epsi$
    then $ephi_1,\dots,\ephi_n\eentails\epsi$
  \end{argument}
  Second, our system is \emph{complete} in that every valid argument is also
  provable.
  \begin{argument}
  \item[ ] if $ephi_1,\dots,\ephi_n\eentails\epsi$
    then $ephi_1,\dots,\ephi_n\eproves\epsi$
  \end{argument}
\end{observation}

Corollary to this observation is the practical equivalence of several analogous
proof-theoretic and model-theoretic concepts.

\begin{definition}[Properties of First-Order Logic]
  Since first-order logic is sound and complete, the following pairs of concepts
  are analogues.\\
  \vspace{0.2cm}
  \begin{center}
  \begin{tabular}{c c c}
    \textbf{model-theoretic} &     & \textbf{proof-theoretic} \\
    % \hline
    valid                    & iff & provable           \\
    validity                 & iff & theorem            \\
    equivalent               & iff & interderivable     \\
    jointly satisfiable      & iff & jointly equivalent \\
  \end{tabular}
  \end{center}
\end{definition}

Strategies for proving these concepts in FOL metalogic can be summarized in the
usual table.

\begin{center}
  \begin{tabular}{c c c}
    & \textbf{yes} & \textbf{no} \\
    \hline
    &&\\    
    \emph{valid (sentence)?}    &
    prove $\efollows \ephi$ &
    find \ephi counter-model \\
    &&\\
    \emph{contradiction?}       &
    prove $\ephi\efollows\ebot$ &
    find $\enot\ephi$ counter-model \\
    &&\\    
    \emph{equivalent?}          &
    \makecell{prove $\ephi\efollows\epsi$ and\\$\epsi\efollows\ephi$} &
    \makecell{find model where \\\ephi and \epsi differ} \\
    &&\\    
    \emph{jointly satisfiable?} &
    \makecell{model where $\ephi_i$\\are all true} &
    prove $\ephi_1\dots\ephi_n\efollows\ebot$ \\
    &&\\    
    \emph{valid (argument)?}    &
    prove $\ephi_1\dots\ephi_n\efollows\epsi$ &
    \makecell{find model with $\ephi_1\dots\ephi_n$\\all true but \epsi false}\\
  \end{tabular}
\end{center}

%% Exercises: quantifier equivalence proofs (pp308-11)

